% ============================================================
% OCR Presentation - Professional Styled Version
% Compile with: xelatex ocr-presentation-styled.tex
% ============================================================
\documentclass[aspectratio=169, 10pt]{beamer}

% ── Font & Language (XeLaTeX) ──────────────────────────────
\usepackage{fontspec}
\setmainfont{Lato}[
 BoldFont = Lato Bold,
 ItalicFont = Lato Italic,
 BoldItalicFont = Lato Bold Italic
]
\setsansfont{Lato}[
 BoldFont = Lato Bold,
 ItalicFont = Lato Italic,
 BoldItalicFont = Lato Bold Italic
]
\setmonofont{Noto Sans Mono}

\usepackage{polyglossia}
\setmainlanguage{vietnamese}

% ── Custom Color Palette ───────────────────────────────────
\definecolor{primary}{HTML}{1B2A4A} % Deep navy
\definecolor{accent}{HTML}{2E86DE} % Vibrant blue
\definecolor{accentlight}{HTML}{D6EAF8} % Light blue bg
\definecolor{accentdark}{HTML}{1A5276} % Dark blue
\definecolor{success}{HTML}{27AE60} % Green
\definecolor{warning}{HTML}{F39C12} % Orange
\definecolor{danger}{HTML}{E74C3C} % Red
\definecolor{graytext}{HTML}{7F8C8D} % Gray text
\definecolor{lightbg}{HTML}{F8F9FA} % Light background
\definecolor{darktext}{HTML}{2C3E50} % Dark text
\definecolor{sectionbg}{HTML}{EBF5FB} % Section background
\definecolor{codegreen}{HTML}{27AE60}
\definecolor{codepurple}{HTML}{8E44AD}
\definecolor{backcolour}{HTML}{F5F6FA}

% ── Beamer Theme Customization ─────────────────────────────
\usetheme{Madrid}
\usecolortheme{whale}

% Override beamercolor
\setbeamercolor{structure}{fg=accent}
\setbeamercolor{frametitle}{fg=white, bg=primary}
\setbeamercolor{title}{fg=white, bg=primary}
\setbeamercolor{subtitle}{fg=accentlight}
\setbeamercolor{author}{fg=white}
\setbeamercolor{institute}{fg=accentlight}
\setbeamercolor{date}{fg=accentlight}
\setbeamercolor{section in toc}{fg=primary}
\setbeamercolor{subsection in toc}{fg=accentdark}
\setbeamercolor{normal text}{fg=darktext}
\setbeamercolor{block title}{fg=white, bg=accent}
\setbeamercolor{block body}{fg=darktext, bg=lightbg}
\setbeamercolor{block alerted title}{fg=white, bg=danger}
\setbeamercolor{block alerted body}{fg=darktext, bg=red!5}
\setbeamercolor{block example title}{fg=white, bg=success}
\setbeamercolor{block example body}{fg=darktext, bg=green!5}
\setbeamercolor{item}{fg=accent}
\setbeamercolor{subitem}{fg=accentdark}
\setbeamercolor{button}{bg=accent, fg=white}

% Custom footline with slide number + logo bar
\setbeamercolor{footline}{fg=white, bg=primary}
\setbeamertemplate{footline}{
 \leavevmode%
 \hbox{%
 \begin{beamercolorbox}[wd=.4\paperwidth,ht=2.5ex,dp=1ex,left]{footline}%
 \hspace*{2ex}\insertshorttitle
 \end{beamercolorbox}%
 \begin{beamercolorbox}[wd=.35\paperwidth,ht=2.5ex,dp=1ex,center]{footline}%
 \insertshortauthor
 \end{beamercolorbox}%
 \begin{beamercolorbox}[wd=.25\paperwidth,ht=2.5ex,dp=1ex,right]{footline}%
 \insertframenumber{} / \inserttotalframenumber\hspace*{2ex}
 \end{beamercolorbox}}%
 \vskip0pt%
}

% Custom headline with accent bar
\setbeamercolor{headline}{fg=white, bg=accent}
\setbeamertemplate{headline}{
 \leavevmode%
 \hbox{%
 \begin{beamercolorbox}[wd=\paperwidth,ht=3pt,dp=0pt]{headline}%
 \end{beamercolorbox}}%
 \vskip0pt%
}

% Remove navigation symbols
\setbeamertemplate{navigation symbols}{}

% Custom title page
\setbeamertemplate{title page}{
 \vbox{}
 \vfill
 \begingroup
 \centering
 \begin{beamercolorbox}[sep=8pt,center,shadow=true,rounded=true]{title}
 \usebeamerfont{title}\inserttitle\par%
 \ifx\insertsubtitle\@empty%
 \else%
 \vskip0.5em%
 {\usebeamerfont{subtitle}\usebeamercolor[fg]{subtitle}\insertsubtitle\par}%
 \fi%
 \end{beamercolorbox}%
 \vskip1em\par
 \begin{beamercolorbox}[sep=8pt,center]{author}
 \usebeamerfont{author}\insertauthor
 \end{beamercolorbox}
 \begin{beamercolorbox}[sep=4pt,center]{institute}
 \usebeamerfont{institute}\insertinstitute
 \end{beamercolorbox}
 \begin{beamercolorbox}[sep=8pt,center]{date}
 \usebeamerfont{date}\insertdate
 \end{beamercolorbox}
 \endgroup
 \vfill
}

% ── Packages ───────────────────────────────────────────────
\usepackage{graphicx}
\usepackage{amsmath,amssymb}
\usepackage{booktabs}
\usepackage{tikz}
\usetikzlibrary{shapes,arrows,positioning,shadows.blur,calc}
\usepackage{listings}
\usepackage{xcolor}
\usepackage{multicol}
\usepackage{tcolorbox}
\tcbuselibrary{skins,breakable}
% \usepackage{fontawesome5} % removed
\usepackage{tabularx}
\usepackage{colortbl}

% ── Code Listing Style ─────────────────────────────────────
\lstdefinestyle{mystyle}{
 backgroundcolor=\color{backcolour},
 commentstyle=\color{codegreen}\itshape,
 keywordstyle=\color{accent}\bfseries,
 numberstyle=\tiny\color{graytext},
 stringstyle=\color{codepurple},
 basicstyle=\ttfamily\footnotesize,
 breaklines=true,
 frame=none,
 xleftmargin=8pt,
 framexleftmargin=8pt,
 rulecolor=\color{accent!30},
 morekeywords={store_name, date, total, items, name, qty, price}
}
\lstset{style=mystyle}

% ── Custom Boxes ───────────────────────────────────────────
\newtcolorbox{infobox}[1][]{
 colback=accentlight,
 colframe=accent,
 coltitle=white,
 fonttitle=\bfseries,
 title=#1,
 rounded corners,
 boxrule=0.8pt,
 left=6pt, right=6pt, top=4pt, bottom=4pt
}

\newtcolorbox{warnbox}[1][]{
 colback=orange!8,
 colframe=warning,
 coltitle=white,
 fonttitle=\bfseries,
 title=#1,
 rounded corners,
 boxrule=0.8pt,
 left=6pt, right=6pt, top=4pt, bottom=4pt
}

\newtcolorbox{successbox}[1][]{
 colback=green!8,
 colframe=success,
 coltitle=white,
 fonttitle=\bfseries,
 title=#1,
 rounded corners,
 boxrule=0.8pt,
 left=6pt, right=6pt, top=4pt, bottom=4pt
}

% ── Metadata ───────────────────────────────────────────────
\title[Trích xuất Hóa đơn VN]{\textbf{Trích xuất Thông tin Hóa đơn Tiếng Việt từ Ảnh}}
\subtitle{Pipeline PaddleOCR -- VietOCR -- LayoutLMv3}
\author{Đồ án Môn học}
\institute{Học viện Công nghệ Bưu chính Viễn thông \\ Khoa Công nghệ Thông tin 2}
\date{Tháng 6, 2026}

% ============================================================
\begin{document}

%---------------------------------------------------------
% Slide 1: Title
%---------------------------------------------------------
\begin{frame}[plain]
 \titlepage
\end{frame}

%---------------------------------------------------------
% Slide 2: Outline
%---------------------------------------------------------
\begin{frame}{Nội dung trình bày}
 \tableofcontents
\end{frame}

%---------------------------------------------------------
% Section 1: MỞ ĐẦU
%---------------------------------------------------------
\section{Mở đầu}

% Slide 3
\begin{frame}{Bối cảnh bài toán}
 \begin{itemize}
 \item[$\bullet$] Mỗi ngày, hàng triệu giao dịch thương mại tại Việt Nam tạo ra hàng triệu tờ hóa đơn giấy
 \item[$\bullet$] Kế toán viên gõ tay từng con số $\rightarrow$ sai sót nhập liệu thường xuyên
 \item[$\bullet$] Máy tính đã có thể đọc chữ in từ lâu, nhưng đọc hóa đơn thực tế vẫn là thách thức:
 \begin{itemize}
 \item Font chữ tùy ý, góc chụp nghiêng
 \item Ánh sáng không đều
 \item Tiếng Việt với đầy đủ dấu thanh
 \end{itemize}
 \item[$\bullet$] Cần biến \textbf{ảnh hóa đơn} thành \textbf{dữ liệu có cấu trúc}
 \end{itemize}
\end{frame}

% Slide 4
\begin{frame}{Các tình huống thực tế}
 \begin{infobox}[Kịch bản A -- Kế toán doanh nghiệp nhỏ]
 300 hóa đơn công tác phí, mất 2--3 ngày nhập thủ công
 \end{infobox}
 \vspace{0.2cm}
 \begin{infobox}[Kịch bản B -- Ứng dụng chi tiêu cá nhân]
 Chụp ảnh hóa đơn $\rightarrow$ tự động điền ngày, siêu thị, tổng tiền
 \end{infobox}
 \vspace{0.2cm}
 \begin{infobox}[Kịch bản C -- Kiểm toán \& tuân thủ thuế]
 Xác minh hàng nghìn hóa đơn VAT tự động
 \end{infobox}
 \vspace{0.3cm}
 \centering
 $\Rightarrow$ Nhu cầu chung: \textbf{\textcolor{accent}{Document Information Extraction (DIE)}}
\end{frame}

%---------------------------------------------------------
% Section 2: TỔNG QUAN
%---------------------------------------------------------
\section{Tổng quan bài toán}

% Slide 5
\begin{frame}{Mô hình hóa bài toán -- Ba tầng xử lý}
 \begin{center}
 \begin{tikzpicture}[
 box/.style={rectangle, draw=accent, rounded corners=8pt, fill=accentlight, text width=11cm, minimum height=1.3cm, align=center, font=\small, drop shadow={shadow xshift=1pt, shadow yshift=-1pt, opacity=0.2}},
 arrow/.style={->, thick, >=stealth, color=accent}
 ]
 \node[box, fill=accent!20] (t3) at (0,3.8) {\textbf{\textcolor{primary}{Tầng 3:}} Trích xuất thực thể có cấu trúc (KIE) -- LayoutLMv3};
 \node[box, fill=accent!12] (t2) at (0,2.0) {\textbf{\textcolor{primary}{Tầng 2:}} Nhận dạng ký tự (OCR Recognition) -- VietOCR};
 \node[box, fill=accentlight] (t1) at (0,0) {\textbf{\textcolor{primary}{Tầng 1:}} Phát hiện vùng chữ (Text Detection) -- PaddleOCR};
 \draw[arrow, line width=1.5pt] (t1) -- node[right, font=\small\bfseries, text=accent]{$\downarrow$} (t2);
 \draw[arrow, line width=1.5pt] (t2) -- node[right, font=\small\bfseries, text=accent]{$\downarrow$} (t3);
 \end{tikzpicture}
 \end{center}
\end{frame}

% Slide 6
\begin{frame}{Đầu vào \& Đầu ra}
 \begin{tcolorbox}[colback=accentlight, colframe=accent, title=Đầu vào, rounded corners]
 Ảnh JPEG/PNG hóa đơn chụp bằng điện thoại (nghiêng, nhòe, ánh sáng không đều)
 \end{tcolorbox}
 \vspace{0.3cm}
 \textbf{Đầu ra:} Cấu trúc JSON có ngữ nghĩa
 \vspace{0.2cm}

 \begin{center}
 \begin{tcolorbox}[colback=green!5, colframe=success, rounded corners, width=13cm]
 \texttt{\footnotesize
 \{ store\_name: Circle K, date: 2024-05-27, total: 87000,\newline
 \ \ \ \ items: [ \{name: Cafe sua da, qty: 2, price: 29000\}, ... ] \}}
 \end{tcolorbox}
 \end{center}
\end{frame}

% Slide 7
\begin{frame}{Pipeline xử lý}
 \begin{center}
 \begin{tikzpicture}[
 node distance=1.6cm,
 process/.style={rectangle, draw=accent, rounded corners=6pt, fill=accent!15, minimum height=1.1cm, minimum width=2.8cm, align=center, font=\footnotesize\bfseries, drop shadow={shadow xshift=0.5pt, shadow yshift=-0.5pt, opacity=0.15}},
 data/.style={rectangle, draw=success, rounded corners=6pt, fill=green!10, minimum height=0.9cm, minimum width=2.8cm, align=center, font=\footnotesize, drop shadow={shadow xshift=0.5pt, shadow yshift=-0.5pt, opacity=0.15}},
 arrow/.style={->, thick, >=stealth, color=accent}
 ]
 \node[data] (img) {Ảnh hóa đơn};
 \node[process, right=of img] (det) {PaddleOCR\\Detection};
 \node[data, right=of det] (bbox) {Bounding\\Boxes};
 \node[process, right=of bbox] (rec) {VietOCR\\Recognition};
 \node[data, right=of rec] (text) {Text + BBox};
 \node[process, below=1.6cm of rec] (kie) {LayoutLMv3\\KIE};
 \node[data, left=of kie] (json) {JSON\\Cấu trúc};

 \draw[arrow] (img) -- (det);
 \draw[arrow] (det) -- (bbox);
 \draw[arrow] (bbox) -- (rec);
 \draw[arrow] (rec) -- (text);
 \draw[arrow] (text) -- (kie);
 \draw[arrow] (kie) -- (json);
 \end{tikzpicture}
 \end{center}
\end{frame}

%---------------------------------------------------------
% Section 3: PHÂN LOẠI & TIẾN HÓA
%---------------------------------------------------------
\section{Phân loại và tiến hóa}

% Slide 8
\begin{frame}{Text Detection -- Từ sliding window đến anchor-free}
 \begin{columns}[T]
 \begin{column}{0.48\textwidth}
 \begin{tcolorbox}[colback=gray!8, colframe=graytext, title={Thế hệ 1 (trước 2015)}, fonttitle=\bfseries, rounded corners, boxrule=0.5pt]
 \begin{itemize}
 \item MSER, SWT
 \item Phân tích pixel theo màu sắc, độ rộng nét bút
 \item Chỉ hoạt động với ảnh sạch, chữ đứng
 \end{itemize}
 \end{tcolorbox}
 \vspace{0.2cm}
 \begin{tcolorbox}[colback=orange!8, colframe=warning, title={Thế hệ 2 (2015--2017)}, fonttitle=\bfseries, rounded corners, boxrule=0.5pt]
 \begin{itemize}
 \item CTPN, TextBoxes++, SegLink
 \item CNN-based với Anchor Boxes
 \item Khó phát hiện văn bản cong
 \end{itemize}
 \end{tcolorbox}
 \end{column}
 \begin{column}{0.48\textwidth}
 \begin{tcolorbox}[colback=accentlight, colframe=accent, title={Thế hệ 3 (2018--nay)}, fonttitle=\bfseries, rounded corners, boxrule=0.8pt]
 \begin{itemize}
 \item \textbf{DB (Differentiable Binarization)}
 \item Xương sống của PaddleOCR
 \item Dự đoán probability map per pixel
 \item Ngưỡng có thể học được (differentiable threshold)
 \item Tối ưu cho CPU/mobile
 \end{itemize}
 \end{tcolorbox}
 \end{column}
 \end{columns}
\end{frame}

% Slide 9
\begin{frame}{Text Recognition -- Từ template matching đến Attention}
 \begin{columns}[T]
 \begin{column}{0.48\textwidth}
 \begin{tcolorbox}[colback=gray!8, colframe=graytext, title={Thế hệ 1 (trước 2014)}, fonttitle=\bfseries, rounded corners, boxrule=0.5pt]
 \begin{itemize}
 \item Tesseract v3, template matching
 \item Tách từng ký tự, so sánh template
 \item Yêu cầu font rõ ràng
 \end{itemize}
 \end{tcolorbox}
 \vspace{0.2cm}
 \begin{tcolorbox}[colback=orange!8, colframe=warning, title={Thế hệ 2 (2015--2018)}, fonttitle=\bfseries, rounded corners, boxrule=0.5pt]
 \begin{itemize}
 \item \textbf{CRNN + CTC}
 \item CNN + BiLSTM + CTC loss
 \item Nền tảng OCR hiện đại
 \end{itemize}
 \end{tcolorbox}
 \end{column}
 \begin{column}{0.48\textwidth}
 \begin{tcolorbox}[colback=accentlight, colframe=accent, title={Thế hệ 3 (2018--nay)}, fonttitle=\bfseries, rounded corners, boxrule=0.8pt]
 \begin{itemize}
 \item \textbf{VietOCR -- Seq2Seq + Attention}
 \item VGG/ResNet encoder + Transformer decoder
 \item Xử lý dấu thanh tiếng Việt tốt hơn
 \item Pretrain trên 5M+ ảnh text tiếng Việt
 \item 60+ font chữ thông dụng
 \end{itemize}
 \end{tcolorbox}
 \end{column}
 \end{columns}
\end{frame}

% Slide 10
\begin{frame}{Key Information Extraction -- Từ rule-based đến Multi-modal}
 \begin{columns}[T]
 \begin{column}{0.48\textwidth}
 \begin{tcolorbox}[colback=gray!8, colframe=graytext, title={Thế hệ 1 -- Rule-based}, fonttitle=\bfseries, rounded corners, boxrule=0.5pt]
 \begin{itemize}
 \item Regex, heuristic
 \item Nhanh, không cần training
 \item Vỡ ngay khi gặp format mới
 \end{itemize}
 \end{tcolorbox}
 \vspace{0.2cm}
 \begin{tcolorbox}[colback=orange!8, colframe=warning, title={Thế hệ 2 -- BERT NER}, fonttitle=\bfseries, rounded corners, boxrule=0.5pt]
 \begin{itemize}
 \item Token classification
 \item Bỏ qua thông tin vị trí không gian
 \end{itemize}
 \end{tcolorbox}
 \end{column}
 \begin{column}{0.48\textwidth}
 \begin{tcolorbox}[colback=accentlight, colframe=accent, title={Thế hệ 3 -- LayoutLMv3}, fonttitle=\bfseries, rounded corners, boxrule=0.8pt]
 \begin{itemize}
 \item Microsoft, 2022
 \item Kết hợp 3 nguồn: text OCR + tọa độ bbox + patch ảnh
 \item Unified self-attention qua cả 3 modality
 \item Pre-training: MLM + MIM + Text-Image Alignment
 \item \textbf{Đột phá} cho Document AI
 \end{itemize}
 \end{tcolorbox}
 \end{column}
 \end{columns}
\end{frame}

%---------------------------------------------------------
% Section 4: TRỰC GIÁC
%---------------------------------------------------------
\section{Trực giác và tư duy}

% Slide 11
\begin{frame}{Con người đọc hóa đơn như thế nào?}
 \begin{enumerate}
 \item[$\bullet$] \textbf{Vị trí không gian mang ngữ nghĩa:}
 \begin{itemize}
 \item ``Tổng cộng'' $\rightarrow$ góc dưới phải, font lớn, in đậm
 \item Tên cửa hàng $\rightarrow$ trên cùng, canh giữa
 \end{itemize}
 \item[$\bullet$] \textbf{Căn lề là thông tin:}
 \begin{itemize}
 \item Giá tiền căn phải, tên mặt hàng căn trái
 \end{itemize}
 \item[$\bullet$] \textbf{Gần nhau $\rightarrow$ liên quan:}
 \begin{itemize}
 \item ``2 x 29,000'' và ``58,000'' cùng dòng $\rightarrow$ cùng mặt hàng
 \end{itemize}
 \item[$\bullet$] \textbf{Pattern lặp lại trong vùng bảng:}
 \begin{itemize}
 \item 5 dòng có cấu trúc [tên $|$ số $|$ số] $\rightarrow$ bảng mặt hàng
 \end{itemize}
 \end{enumerate}
\end{frame}

% Slide 12
\begin{frame}{Mô hình mô phỏng trực giác con người}
 \begin{center}
 \begin{tabularx}{\textwidth}{>{\bfseries\color{primary}}p{3.5cm}X}
 \toprule
 \textcolor{primary}{Trực giác} & \textcolor{primary}{Mô hình mô phỏng} \\
 \midrule
 \rowcolor{accentlight}
 Vùng chữ có cấu trúc pixel đặc biệt & PaddleOCR Detection: học probability map \\
 Đọc ký tự cần ngữ cảnh xung quanh & VietOCR: attention ``nhìn lại'' toàn bộ ảnh \\
 \rowcolor{accentlight}
 Vị trí + nội dung + hình ảnh cùng quyết định & LayoutLMv3: fuse 3 luồng trong transformer \\
 \bottomrule
 \end{tabularx}
 \end{center}
 \vspace{0.4cm}
 \begin{warnbox}[Giả định ngầm]
 Mỗi mô hình đều có giả định riêng về dữ liệu input $\rightarrow$ khi giả định vi phạm, mô hình suy giảm chất lượng.
 \end{warnbox}
\end{frame}

%---------------------------------------------------------
% Section 5: TOÁN HỌC HÓA
%---------------------------------------------------------
\section{Toán học hóa}

% Slide 13
\begin{frame}{PaddleOCR -- DB (Differentiable Binarization)}
 \begin{infobox}[Phép toán chủ đạo: Segmentation + Adaptive Thresholding]
 \end{infobox}
 \vspace{0.3cm}

 Mô hình học 2 map song song từ feature map $F$:
 \begin{align*}
 P &= \sigma(\text{Conv}(F)) && \leftarrow \text{\textcolor{accent}{Probability map}} \\
 T &= \sigma(\text{Conv}'(F)) && \leftarrow \text{\textcolor{warning}{Threshold map}} \\
 \hat{B} &= \sigma\left(\frac{P - T}{k}\right) && \leftarrow \text{\textcolor{success}{Binary map (differentiable)}}
 \end{align*}
 \vspace{0.2cm}

 \textbf{Bí quyết:} $\sigma\left(\frac{P-T}{k}\right)$ với $k=50$ xấp xỉ hàm step nhưng vẫn có gradient để backprop.

 \vspace{0.2cm}
 \begin{tcolorbox}[colback=lightbg, colframe=primary, boxrule=0.5pt, rounded corners]
 \textbf{Loss:} $L = L_{\text{bce}}(P, G_p) + \alpha \cdot L_{\text{bce}}(\hat{B}, G_b) + \beta \cdot L_1(T, G_t)$ \\
 \textbf{Tham số:} ResNet-18/50 + FPN, khoảng 10--27M tham số.
 \end{tcolorbox}
\end{frame}

% Slide 14
\begin{frame}{VietOCR -- Seq2Seq với Attention}
 \begin{infobox}[Phép toán chủ đạo: Cross-Attention trong Transformer decoder]
 \end{infobox}
 \vspace{0.3cm}

 \begin{align*}
 F &= \text{CNN\_Encoder}(\text{ảnh}) \quad \text{shape: } (H/4, W/4, C) \\
 F_{\text{flat}} &= \text{reshape}(F) + \text{positional\_embed} \quad \text{shape: } (S, C) \\
 h_t &= \text{TransformerDecoder}(h_{t-1}, F_{\text{flat}}) \\
 p(y_t) &= \text{Softmax}(\text{Linear}(h_t))
 \end{align*}
 \vspace{0.2cm}

 \textbf{Attention:} $\alpha_{t,s} = \text{softmax}\left(\frac{Q_t \cdot K_s}{\sqrt{d_k}}\right)$ -- cho phép ``nhìn'' vào vùng feature tương ứng.

 \vspace{0.2cm}
 \begin{successbox}[Đặc biệt tiếng Việt]
 Khi sinh ``ề'', decoder attend mạnh vào vùng có dấu phụ kép (mũ + huyền). \\
 \textbf{Vocab:} 130--150 ký tự Unicode tiếng Việt + dấu câu + số.
 \end{successbox}
\end{frame}

% Slide 15
\begin{frame}{LayoutLMv3 -- Multi-modal Transformer}
 \begin{infobox}[Phép toán chủ đạo: Unified Self-Attention qua cả 3 modality]
 \end{infobox}
 \vspace{0.3cm}

 \textbf{Input token $i$:}
 \[e_i = E_{\text{text}}(\text{token}_i) + E_{\text{1D}}(\text{pos}_i) + E_{\text{2D}}(x_1, y_1, x_2, y_2, w, h)\]

 \textbf{Image patch $j$:}
 \[v_j = \text{LinearProj}(\text{patch}_j) + E_{\text{img\_pos}}(j)\]

 \textbf{Concatenate $\rightarrow$ Transformer $\rightarrow$ Contextualized representations}
 \vspace{0.3cm}

 \textbf{Pre-training objectives:}
 \begin{itemize}
 \item[$\bullet$] Masked Language Modeling (MLM) trên text
 \item[$\bullet$] Masked Image Modeling (MIM) trên ảnh
 \item[$\bullet$] Text-Image Alignment
 \end{itemize}
\end{frame}

%---------------------------------------------------------
% Section 6: KHOẢNG CÁCH
%---------------------------------------------------------
\section{Khoảng cách giữa lý thuyết và thực tế}

% Slide 16
\begin{frame}{Hàm mất mát và lý do lựa chọn}
 \begin{center}
 \begin{tabularx}{\textwidth}{>{\bfseries}p{2.5cm}p{4cm}X}
 \toprule
 \textcolor{primary}{Tầng} & \textcolor{primary}{Loss Function} & \textcolor{primary}{Tại sao?} \\
 \midrule
 \rowcolor{accentlight}
 Detection & BCE + L1 & Gradient ổn định, robust với outlier \\
 Recognition & Cross-Entropy chuỗi & Smooth, vi phân được \\
 \rowcolor{accentlight}
 KIE & Cross-Entropy/token & Tính được tại từng token \\
 \bottomrule
 \end{tabularx}
 \end{center}
 \vspace{0.4cm}

 \textbf{Tại sao không dùng Metric làm Loss?}
 \begin{itemize}
 \item[\textcolor{danger}{$\bullet$}] Field Exact Match: không có gradient
 \item[\textcolor{danger}{$\bullet$}] CER (edit distance): không liên tục
 \item[\textcolor{danger}{$\bullet$}] Entity F1: phụ thuộc ngưỡng confidence
 \end{itemize}
 \vspace{0.2cm}
 \centering
 $\Rightarrow$ Dùng CE làm Loss, F1/CER/Exact Match chỉ để đánh giá.
\end{frame}

% Slide 17
\begin{frame}{Hệ quả thực tế}
 \begin{alertblock}{Điểm đau lớn nhất}
 Mô hình có thể đạt \textbf{CER rất thấp (98\% ký tự đúng)} nhưng \textbf{Field Exact Match chỉ 70\%} -- vì chỉ cần sai 1 ký tự trong ``87.000'' thành ``87.500'' là cả trường total sai hoàn toàn.
 \end{alertblock}
 \vspace{0.3cm}

 \textbf{Khả năng mở rộng Loss:}
 \begin{itemize}
 \item[\textcolor{success}{$\bullet$}] Thêm class mới: DISCOUNT, TAX\_CODE
 \item[\textcolor{success}{$\bullet$}] Weighted loss: ưu tiên TOTAL, DATE
 \item[\textcolor{success}{$\bullet$}] Consistency constraint: $\sum$(subtotals) $\approx$ total
 \end{itemize}
\end{frame}

%---------------------------------------------------------
% Section 7: THÁCH THỨC
%---------------------------------------------------------
\section{Thách thức và rào cản}

% Slide 18
\begin{frame}{Trở ngại đã giải quyết \& còn tồn đọng}
 \begin{successbox}[Đã giải quyết]
 \begin{itemize}
 \item Văn bản đa hướng, cong $\rightarrow$ DB, ABCNet (bezier curves)
 \item Font chữ đa dạng tiếng Việt $\rightarrow$ VietOCR pretrained + augmentation
 \item Hiểu layout cơ bản $\rightarrow$ 2D positional embedding (LayoutLM)
 \end{itemize}
 \end{successbox}
 \vspace{0.3cm}

 \begin{warnbox}[Còn tồn đọng]
 \begin{itemize}
 \item[$\times$] Ảnh chất lượng thấp (run, nhòe, phản chiếu) -- CER tụt 85--90\%
 \item[$\times$] Hóa đơn in nhiệt bị phai -- contrast rất thấp
 \item[$\times$] Zero-shot generalization -- cần fine-tune cho mỗi template mới
 \item[$\times$] Dữ liệu tiếng Việt khan hiếm, hạ tầng GPU đắt
 \item[$\times$] Relation Extraction phức tạp (20+ mặt hàng)
 \end{itemize}
 \end{warnbox}
\end{frame}

% Slide 19
\begin{frame}{Hướng đi tương lai}
 \begin{enumerate}
 \item[$\bullet$] \textbf{End-to-end Document Understanding:}
 \begin{itemize}
 \item Donut, Nougat -- bỏ qua OCR, ảnh $\rightarrow$ JSON trực tiếp
 \end{itemize}
 \item[$\bullet$] \textbf{Large Vision-Language Models (LVLMs):}
 \begin{itemize}
 \item GPT-4V, Gemini -- zero-shot, không cần fine-tune
 \end{itemize}
 \item[$\bullet$] \textbf{Self-supervised pretraining trên tài liệu Việt:}
 \begin{itemize}
 \item Corpus hóa đơn/biểu mẫu tiếng Việt quy mô lớn
 \end{itemize}
 \item[$\bullet$] \textbf{Synthetic data generation:}
 \begin{itemize}
 \item Tạo hóa đơn giả + domain randomization
 \end{itemize}
 \end{enumerate}
\end{frame}

%---------------------------------------------------------
% Section 8: KẾT LUẬN
%---------------------------------------------------------
\section{Kết luận}

% Slide 20
\begin{frame}{Kết luận}
 \begin{itemize}
 \item[$\bullet$] Bài toán trích xuất hóa đơn tiếng Việt đòi hỏi giải quyết đồng thời 3 thách thức: \textbf{phát hiện vùng chữ}, \textbf{nhận dạng ký tự}, \textbf{hiểu cấu trúc ngữ nghĩa}
 \vspace{0.2cm}
 \item[$\bullet$] Pipeline \textbf{\textcolor{accent}{PaddleOCR $\rightarrow$ VietOCR $\rightarrow$ LayoutLMv3}} là lựa chọn hợp lý hiện tại với pretrained weights chất lượng tốt
 \vspace{0.2cm}
 \item[$\bullet$] Nhược điểm cố hữu: \textbf{\textcolor{danger}{lỗi tích lũy qua từng tầng}}
 \vspace{0.2cm}
 \item[$\bullet$] Ranh giới giữa OCR+NLP pipeline và end-to-end LVLM đang dần mờ đi
 \end{itemize}
\end{frame}

%---------------------------------------------------------
% Slide 21: Cảm ơn
%---------------------------------------------------------
\begin{frame}[plain]
 \begin{tikzpicture}[remember picture, overlay]
 \fill[primary] (current page.south west) rectangle (current page.north east);
 \fill[accent!30] ([yshift=-3cm]current page.north) circle (8cm);
 \end{tikzpicture}
 \vspace{2cm}
 \begin{center}
 {\Huge\bfseries\color{white} Cảm ơn Thầy và Các bạn!}
 \vspace{1cm}

 {\Large\color{accentlight} Xin cảm ơn!}
 \vspace{0.8cm}

 {\normalsize\color{accentlight} Học viện Công nghệ Bưu chính Viễn thông \\ Khoa Công nghệ Thông tin 2}
 \end{center}
\end{frame}

\end{document}
